\documentclass[uplatex,dvipdfmx,a4paper,11pt]{jsarticle}

%url
\usepackage{url}
\usepackage[dvipdfmx]{hyperref}

% --- 余白 ---
\usepackage[top=25mm,bottom=25mm,left=25mm,right=25mm]{geometry}

% --- 表・箇条書き・図 ---
\usepackage{array}
\usepackage{tabularx}
\usepackage{booktabs}
\usepackage{enumitem}
\usepackage{graphicx}
\usepackage{longtable}

% --- 段落間隔（Wordのような見た目に寄せる）---
\setlength{\parindent}{0pt}
\setlength{\parskip}{0.6em}

% --- セクション見出しを「番号. 太字」+ 章末に水平線を入れる ---
\usepackage{titlesec}
\titleformat{\section}
  {\normalfont\Large\bfseries}{\thesection.}{0.6em}{}
\titleformat{\subsection}
  {\normalfont\large\bfseries}{\thesubsection}{0.6em}{}
\titlespacing*{\section}{0pt}{1.2em}{0.6em}
\titlespacing*{\subsection}{0pt}{1.0em}{0.4em}

% --- 章末水平線（手で \sectionrule と書いて区切る）---
\newcommand{\sectionrule}{\par\medskip\noindent\rule{\linewidth}{0.6pt}\par\medskip}

% --- 箇条書きの記号を Word 風（・→○→数字）に揃える ---
\renewcommand{\labelitemi}{\textbullet}
\renewcommand{\labelitemii}{\textopenbullet}
\setlist[itemize]{leftmargin=1.6em,itemsep=0.1em,topsep=0.2em}
\setlist[enumerate]{leftmargin=2.0em,itemsep=0.1em,topsep=0.2em}

% \textopenbullet が無い環境向けのフォールバック（jsarticle標準では○が出る）
\providecommand{\textopenbullet}{\ensuremath{\circ}}


\begin{document}


{\LARGE\bfseries 作業ログ}\par\bigskip

報告者：佐藤涼菜

報告日：5月27日

プロジェクト名：マジカル☆観覧車

\sectionrule


\section{進捗サマリー（要約）}

\begin{itemize}
  \item 全体進捗率：２０[％]
  \item マイルストーン達成状況：おおよそ達成
  \item 主要成果：
    \begin{itemize}
      \item {計画書の変更}
      \item {LEDマトリクスの準備および動作確認}
      \item {指揮デバイスの製図}
      \item {LEDMatrixの制御}
    \end{itemize}
  \item 重大課題（影響度：高）：
    \begin{itemize}
      \item {特になし．}
    \end{itemize}
\end{itemize}

\sectionrule


\section{作業ログ（詳細トラッキング）}

\subsection{作業ログ}

\renewcommand{\arraystretch}{1.2}

\small

\begin{longtable}{|p{3cm}|c|c|c|c|p{3.5cm}|p{4cm}|}
\caption{作業ログ} \\
\hline
作業項目 & 開始日 & 予定完了日 & 状態 & 進捗率 & 依存関係・ブロッカー & コメント \\
\hline
\endfirsthead

\hline
作業項目 & 開始日 & 予定完了日 & 状態 & 進捗率 & 依存関係・ブロッカー & コメント \\
\hline
\endhead

\hline
\multicolumn{7}{r}{次ページへ続く} \\
\hline
\endfoot

\hline
\endlastfoot

計画書の図の差し替え &
5/20 &
5/21 &
完了 &
100\% &
無 &
特になし． \\
\hline

LEDMatrixの動作確認 &
5/22 &
5/22 &
完了 &
100\% &
無 &
特になし． \\
\hline

作業ログの記入 &
5/26 &
5/26 &
完了 &
100\% &
無 &
特になし． \\
\hline

指揮デバイスの製図 &
5/24 &
5/28 &
進行中 &
35\% &
無 &
特になし． \\
\hline

LEDMatrixの制御 &
5/23 &
5/30 &
進行中 &
40\% &
無 &
特になし． \\
\hline



\end{longtable}


\subsection{日ごとの作業ログ}
 

    \renewcommand{\arraystretch}{1.2}

    \small
    
    \begin{longtable}{|c|p{3cm}|c|c|c|c|p{5cm}|}
    \caption{日ごとの作業ログ} \\
    \hline
    日付 & 作業項目 & 時間 & 開始時の状態 & 終了時の状態 & 進捗率 & コメント \\
    \hline
    \endfirsthead
    
    \hline
    日付 & 作業項目 & 時間 & 開始時の状態 & 終了時の状態 & 進捗率 & コメント \\
    \hline
    \endhead
    
    \hline
    \multicolumn{7}{r}{次ページへ続く} \\
    \hline
    \endfoot
    
    \hline
    \endlastfoot
    
    5/20 & 計画書の図変更 & 2h &
    未着手 &
    進行中 &
    30\% &
    発表において発生した変更に該当する担当箇所の図の差し替え作業． \\
    \hline
    
    5/21 & 計画書の図変更 & 2h &
    進行中 &
    完了 &
    100\% &
    特になし． \\
    \hline
    
    5/22 & 役割分担の確認および全体の計画書の修正 & 1h &
    未着手 &
    完了 &
    100\% &
    計画書変更に伴い実装役割分担の再検討を行った．
    全員で誤字脱字や不備の修正を行った． \\
    \hline
    
    5/22 & 部品の用意 & 30min &
    未着手 &
    完了 &
    100\% &
    Arduinoに搭載されているLEDMatrixの動作確認，
    および指揮デバイスに使用する部品の確認を行った． \\
    \hline
    
    5/23 & LEDMatrixの制御 & 3h &
    未着手 &
    進行中 &
    20\% &
    部品用意が滞っていたため，予定よりも早めに着手した． \\
    \hline
    
    5/24 & LEDMatrixの制御 & 2.5h &
    進行中 &
    進行中 &
    30\% &
    特になし． \\
    \hline
    
    〃 & 指揮デバイスの製図 & 3h &
    未着手 &
    進行中 &
    25\% &
    Tinkercadの3Dデザインにて製図を作成． \\

    \hline

    5/25 & 指揮デバイスの製図 & 1.5h &
    進行中 &
    進行中 &
    30\% &
    特になし． \\
    \hline

    5/26 & 作業ログの記入 & 2h &
    未着手 &
    完了 &
    100\% &
    特になし． \\  

    \hline

    〃 & LEDMatrixの制御 & 2h &
    進行中 &
    進行中 &
    40\% &
    特になし． \\ 

    \hline
    
    〃 & 指揮デバイスの製図 & 1h &
    進行中 &
    進行中 &
    35\% &
    特になし． \\ 
    
    \end{longtable}

\sectionrule

\section{計画時のGCと現在のGC}
最終計画書に記載されたガントチャート(GC)を図\ref{fig:gant1}に示す．

\begin{figure}[htbp] % h:ここ, t:上, b:下, p:独立ページ
    \centering % 中央揃え
    \includegraphics[width=0.8\columnwidth]{gant1.png} % 横幅を列幅の80%に指定
    \caption{計画時のガントチャート} % 図の説明
    \label{fig:gant1} % 本文中で参照するためのラベル
\end{figure}

次に，現在のGCについて図\ref{fig:gant2}に示す．達成した項目は緑色で塗られている．
担当箇所はLEDマトリクスの用意(No.2121)，本体・センサの実装(No.2112),そしてLEDマトリクスの制御(No.2231)である．
No.2112とNo.2231について，作業を前倒ししている．また，27日に部品を入手でき次第，指揮デバイスの実装に着手する．

\begin{figure}[htbp] % h:ここ, t:上, b:下, p:独立ページ
    \centering % 中央揃え
    \includegraphics[width=0.8\columnwidth]{gant2.png} % 横幅を列幅の80%に指定
    \caption{現在のガントチャート} % 図の説明
    \label{fig:gant2} % 本文中で参照するためのラベル
\end{figure}


\sectionrule


\section{課題管理}

\begin{tabularx}{\linewidth}{@{}lXllXll@{}}
\toprule
課題 & 発生原因 & 影響度 & 優先度 & 対策 & 期限 & 状態 \\
\midrule
進捗が芳しくない & 予備知識が少なく，調べながらの実装であるから. & 低 & 低 & 空いた時間を有効的に活用し，まとまった時間に実装が進むようにする & - & 対応中 \\
\bottomrule
\end{tabularx}

\medskip
※影響度＝プロジェクトへのインパクト

※優先度＝対応の緊急性

\sectionrule


\section{AI利用状況と検証（再現性重視）}

\subsection{利用内容}

\begin{itemize}
  \item 利用目的：設計・実装支援・レビュー

  \item 入力（プロンプト要旨）：
    \begin{itemize}
      \item 指揮デバイスの3D製図を作成する方法や，
      Tinkercadを用いたモデリング手順について質問した．
      
      \item Arduinoに搭載されているLEDMatrixの制御方法，
      表示プログラムの作成方法，
      およびライブラリの利用方法について質問した．
     
    \end{itemize}

  \item 出力概要：
    \begin{itemize}
      \item Tinkercadにおける3Dデザイン機能の利用方法や，
      指揮デバイスの製図手順に関する説明が得られた．
      
      \item ArduinoでLEDMatrixを制御するための
      サンプルコードやライブラリ利用例が得られた．
      
    
    \end{itemize}
\end{itemize}


\sectionrule

\subsection{検証方法}

\begin{itemize}

    \item 検証手法：
      \begin{itemize}
        \item AIが生成した製図手順を用いて，
        Tinkercad上で実際に指揮デバイスの3Dモデルを作成し，
        意図した形状を再現可能か確認した．
  
        \item AIが生成したLEDMatrix制御プログラムを
        Arduino IDE上でコンパイルおよび実行し，
        LEDMatrixへの文字表示や点灯制御が正常に行われるか確認した．
  
        \item AIの出力内容とArduinoおよびTinkercadの公式ドキュメントを比較し，
        使用関数や実装方法に誤りがないか確認した．
      \end{itemize}

      \item 参照資料：
      \begin{itemize}
        \item Arduino公式ドキュメント：
        \url{https://docs.arduino.cc/}
  
        \item Arduino Tutorials：
        \url{https://docs.arduino.cc/tutorials/}
  
        \item Tinkercad公式サイト：
        \url{https://www.tinkercad.com/}
  
        \item Tinkercad Learn：
        \url{https://www.tinkercad.com/learn/}
      \end{itemize}

      \item 検証手順：
    \begin{enumerate}
      \item AIへ指揮デバイスの製図方法，
      LEDMatrix制御方法，
      およびArduinoライブラリの使用方法について質問した．

      \item AIが生成した内容をもとに，
      Tinkercad上で3Dモデルおよび配線図を作成した．

      \item AIが生成したサンプルコードをArduino IDEへ入力し，
      コンパイルエラーの有無を確認した．

      \item Arduinoへプログラムを書き込み，
      LEDMatrixへの文字表示および点灯制御が正常に動作するか確認した．

      \item AIの出力内容をArduinoおよびTinkercadの公式ドキュメントと比較し，
      実装内容に誤りがないことを確認した．
    \end{enumerate}

\end{itemize}



\sectionrule

\subsection{ファクトチェック結果}


\begin{itemize}

    \item 一致点：
      \begin{itemize}
        \item AIが提示したArduinoライブラリの利用方法，
        LEDMatrix制御関数，
        および基本的なプログラム構成は，
        Arduino公式ドキュメントの内容と一致していた．
  
        \item AIが提示したTinkercadでの製図手順についても，
        公式チュートリアルに記載されている操作方法と一致していた．
  
        \item LEDMatrixへの文字表示や点灯制御について，
        実際にArduino上で正常動作することを確認した．
      \end{itemize}

      \item 不一致・疑義：
    \begin{itemize}
      \item 一部のサンプルコードにおいて，
      使用しているライブラリ名や関数名が実際の環境と異なる場合があり，
      修正が必要であった．

      \item Tinkercad上での部品配置について，
      実際の使用環境ではサイズ調整が必要な箇所が存在した．
    \end{itemize}

  \item 最終判断：
    \begin{itemize}
      \item 一部採用
    \end{itemize}

    \item 根拠：
    \begin{itemize}
      \item Arduino公式ドキュメントおよび
      Tinkercad公式チュートリアルとの比較により，
      AIの出力内容がおおむね正確であることを確認したため．

      \item 実際にLEDMatrixの制御および
      指揮デバイスの製図を行い，
      正常に動作・作成できたため．
    \end{itemize}

\end{itemize}






\sectionrule


\section{次回計画}


\begin{itemize}

  \item 目標：指揮デバイスの完成，LEDMatrixの制御
    \begin{itemize}
      \item 指揮デバイスの3Dモデル設計を完了し，
      製図を基に実際に作成する．

      \item Arduinoを用いたLEDMatrixの制御プログラムを完成させ，
      文字表示を実装する．

    \end{itemize}

  \item 達成基準：
    \begin{itemize}
      \item 指揮デバイスの製図データおよび配線図が完成していること．
      \item 指揮デバイスが完成していること．

      \item LEDMatrixに対して指定した文字列や表示パターンを
      正常に表示できること．

  
    \end{itemize}

  \item 期限：
    \begin{itemize}
      \item 2026/6/3
    \end{itemize}

\end{itemize}


\sectionrule


\section{その他}

\begin{itemize}
  \item {}特になし．
\end{itemize}

\sectionrule


\section{付録}

\begin{itemize}
  \item 添付資料：
    \begin{itemize}
      \item ソースコード：\url{https://github.com/Sato907/3S_HCK1/tree/main/HCK1_group/HCK1_group_mywork/Display}
    %  \item 回路図／設計図：[ファイル]
    \end{itemize}

\end{itemize}

\sectionrule

\end{document}
